\documentclass[10pt,a4paper]{amsart}
\usepackage[margin=19mm]{geometry}
\usepackage{amsmath,amssymb,mathtools}
\usepackage[scr=rsfs,cal=euler]{mathalfa}
\usepackage{xfrac}
\usepackage[unicode,hidelinks,bookmarks=false]{hyperref}
\newcommand{\ie}{i.e.\ }
\newcommand{\cf}{cf.\ }
\newcommand{\resp}{respectively\ }
\newcommand{\Subsection}{\subsection}
\providecommand{\nobreakdash}{\nobreak\hbox{-}}
\setlength{\emergencystretch}{2em}
\multlinegap0pt
\usepackage{bm}
\usepackage{bbm}
\usepackage{dsfont}
\usepackage{xspace}
\usepackage{cancel}
\usepackage{mathtools}
\usepackage[shortlabels]{enumitem}
\usepackage{amsthm}
\makeatletter
\@ifundefined{assumption}{\newtheorem{assumption}{Assumption}}{}
\@ifundefined{definition}{\newtheorem{definition}{Definition}}{}
\@ifundefined{theorem}{\newtheorem{theorem}{Theorem}}{}
\makeatother
\makeatletter
\newcommand{\Cn}{\text{Cn}}
\expandafter\ifx\csname proof\endcsname\relax
\renewenvironment{proof}[1][\proofname]{\par
  \pushQED{\qed}%
  \normalfont \topsep6\p@\@plus6\p@\relax
  \trivlist
  \item[\hskip\labelsep
        \itshape
    #1\@addpunct{.}]\ignorespaces
}{%
  \popQED\endtrivlist\@endpefalse
}
\fi
\makeatother
\title[Proof-Valid Caching under Premise Erasures: Local Structural]{Proof-Valid Caching under Premise Erasures: Local Structural Limits and Shared-Workload Gains}
\author{Jianfeng Xu}
\date{Source: arXiv:2608.11782v1 (CC BY 4.0)}
\begin{document}
\maketitle

\noindent\emph{Source and licence.} Proof-Valid Caching under Premise Erasures: Local Structural Limits and Shared-Workload Gains. Version arXiv:2608.11782v1 (CC BY 4.0), distributed under the Creative Commons Attribution 4.0 International licence, as recorded in the arXiv licence metadata for this exact version and shown on the abstract page https://arxiv.org/abs/2608.11782v1. Full rights notice: licenses/arxiv-2608.11782-CC-BY-4.0.txt.\par\smallskip

\noindent\textit{Editorial scope.} Counted unit: the Theorem whose source label is \texttt{thm:envelope-rigidity} in the article (source0.tex lines 1682--1739, source label \texttt{thm:envelope-rigidity}), delivered together with the article's own complete original proof of it (source0.tex lines 1741--1805, one \texttt{proof} environment of 65 lines). No mathematics is written, repaired or re-derived: statement and proof are the article's own text line for line. Required context retained uncounted: assump:det-local (lines 447--490); eq:base-irreducibility (lines 480--485); eq:deterministic-local-law (lines 480--485); assump:hereditary-dag (lines 492--536); def:query-local-residual (lines 982--1017); thm:exact-residual-law (lines 1094--1119); eq:exact-residual-probability (lines 1110--1118); assump:and-or-local (lines 1573--1621); eq:and-or-local-law (lines 1591--1597); eq:well-founded-rank (lines 1615--1620). Every definition, display and label of those lines is retained unchanged. Omitted from this package and not redistributed: 1--446, 486--491, 537--981, 1018--1093, 1119--1572, 1598--1614, 1621--1681, 1740--1740, 1806--5051. Nothing inside a retained excerpt is omitted or altered: every retained body is delivered exactly as the article's lines are written, so no omission receipt of any kind accompanies this package. Citations: the retained excerpts carry no citation command at all, so no bibliography entry of the article is transcribed and no reference is redistributed. Reference adaptation: none was needed and none was made; every reference inside the delivered unit resolves inside it, so no unresolved reference or citation is produced. Printed slips kept literal: no printed word, symbol, hypothesis or number of the counted statement or of its proof was altered. Layout and class: the article's own class is \texttt{IEEEtran} with packages including amsmath, amsfonts, amssymb, amsthm, mathtools, mathrsfs, algorithm, algpseudocode, array, subfig, textcomp, stfloats, url, multirow; the wrapper is the lane's pinned \texttt{amsart} editorial wrapper at 10pt on A4, which restates the theorem-like environments the retained text needs and adds the article's own macro definitions that the retained text needs (\texttt{\textbackslash Cn}), with extra wrapper packages (bm, bbm, dsfont, xspace, cancel, mathtools, enumitem). The delivered numbering is the unit's own. Assets: no retained span contains a figure environment, a graphics-inclusion command, a PDF-page inclusion, a table or a TikZ picture; no image file is copied into this package, the package declares no asset dependency and no image file is redistributed with it. Licence: version arXiv:2608.11782v1 of this article is distributed under the Creative Commons Attribution 4.0 International licence, as recorded in the arXiv licence metadata for this exact version and shown on the abstract page https://arxiv.org/abs/2608.11782v1. A full rights notice is delivered at licenses/arxiv-2608.11782-CC-BY-4.0.txt. Characters: every retained excerpt is pure ASCII, so the delivered engine, XeLaTeX with its default Latin Modern text font, reports no missing character.
\begin{assumption}[Effective deterministic local semantics]
\label{assump:det-local}
The operator \(\Cn\) is a finitary, extensive, idempotent, and monotone
closure operator on finite sets, and membership in \(\Cn(\Gamma)\) is
decidable (effectiveness).  There exist a finite relevant region
\[
  \mathcal V_{\mathrm{rel}}(B)\subseteq\Cn(B),
  \qquad
  B\cup\mathcal Q\subseteq\mathcal V_{\mathrm{rel}}(B),
\]
and a computable designated-parent map
\[
  \operatorname{par}:
  \mathcal V_{\mathrm{rel}}(B)\setminus B
  \longrightarrow
  \mathcal V_{\mathrm{rel}}(B)^{<\infty},
\]
where \(\mathcal V_{\mathrm{rel}}(B)^{<\infty}\) denotes the set of finite
tuples over \(\mathcal V_{\mathrm{rel}}(B)\).  The parent relation is
well-founded: backward expansion along \(\operatorname{par}\) from any
\(v\in\mathcal V_{\mathrm{rel}}(B)\) terminates at vertices in \(B\) after
finitely many steps, so the local law \eqref{eq:deterministic-local-law}
below is an unambiguous recursion rather than a fixpoint condition.
For every
\(v\in\mathcal V_{\mathrm{rel}}(B)\setminus B\), write
\[
  \operatorname{par}(v)
  =
  (u_1,\ldots,u_{a(v)}),
  \qquad a(v)\ge 1,
\]
where the entries are pairwise distinct.  For every finite
\(\Gamma\subseteq\Cn(B)\), the following laws hold:
\begin{align}
  b\in\Cn(\Gamma) &\Longleftrightarrow b\in\Gamma, \qquad b\in B, \label{eq:base-irreducibility}\\
  v\in\Cn(\Gamma) &\Longleftrightarrow 
    \bigl[v\in\Gamma \text{ or } \{u_1,\ldots,u_{a(v)}\}\subseteq\Cn(\Gamma)\bigr], \notag\\
  &\qquad\qquad v\in\mathcal V_{\mathrm{rel}}(B)\setminus B. \label{eq:deterministic-local-law}
\end{align}
Thus a relevant derived object is available either because it is stored
explicitly or because all members of its designated parent tuple are
derivable, whereas a base premise is available only when it is explicitly
present.
\end{assumption}

\begin{assumption}[Finite leaf-hereditary canonical derivation DAG]
\label{assump:hereditary-dag}
For every \(q\in\mathcal Q\), repeatedly expanding from \(q\) along the
designated parent map \(\operatorname{par}(\cdot)\) terminates after finitely
many steps at vertices in \(B\) and yields a finite acyclic graph
\[
  G(q,B)=\bigl(V(G(q,B)),E(G(q,B))\bigr).
\]
Edges are oriented from prerequisites to consequences, so \(q\) is the
distinguished sink and
\[
  V_0(G(q,B))
  :=
  V(G(q,B))\cap B
\]
is the set of source vertices in the edge orientation and terminal leaves
under backward expansion.

For every non-leaf vertex
\(v\in V(G(q,B))\setminus B\), if
\[
  \operatorname{par}(v)=(u_1,\ldots,u_{a(v)}),
\]
then the in-neighborhood of \(v\) in \(G(q,B)\) is exactly the underlying set
\[
  \{u_1,\ldots,u_{a(v)}\}.
\]
Whenever graph-theoretic membership or adjacency is considered, the ordered
tuple \(\operatorname{par}(v)\) is identified with its underlying set; its
ordering is retained only as part of the canonical description.

For every \(v\in V(G(q,B))\), let \(G_q[v]\) denote the predecessor-induced
sub-DAG of \(G(q,B)\) rooted backward at \(v\), and let \(G(v,B)\) denote the
DAG obtained by restarting the same canonical expansion from \(v\).  Then
\begin{equation}
  G_q[v]=G(v,B).
\label{eq:hereditary-subdag}
\end{equation}
When \(v\in B\), both graphs are understood as the singleton graph on \(v\).
In particular,
\begin{equation}
  V_0(G_q[v])=V_0(G(v,B)).
\label{eq:hereditary-leaves}
\end{equation}
\end{assumption}

\begin{definition}[Query-local projection and exposed leaves]
\label{def:query-local-residual}
For a finite transparent cache \(S\subseteq\Cn(B)\), define its projection onto
the canonical query DAG by
\begin{equation}
  U(q,S):=S\cap V(G_q).
\label{eq:query-local-projection}
\end{equation}
For each \(b\in V_0(G_q)\), let \(\mathcal P_q(b)\) denote the set of all
directed paths from \(b\) to \(q\) in \(G_q\), each path being identified
with its vertex set, endpoints included.

A path \(P\in\mathcal P_q(b)\) is \emph{cache-free} if
\[
  V(P)\cap U(q,S)=\varnothing.
\]
The residual exposed-leaf set is
\begin{equation}
  \begin{split}
    D_{\mathrm{exp}}(q,&S)
    := \bigl\{ b\in V_0(G_q) :  \\
    & \exists P\in\mathcal P_q(b) \text{ such that }
      V(P)\cap U(q,S)=\varnothing \bigr\}.
  \end{split}
\label{eq:def-exposed-leaves}
\end{equation}
Its complement
\begin{equation}
  D_{\mathrm{blk}}(q,S)
  :=
  V_0(G_q)\setminus D_{\mathrm{exp}}(q,S)
\label{eq:def-blocked-leaves}
\end{equation}
is the set of base leaves whose every path to \(q\) intersects the query-local
cache.
\end{definition}

\begin{theorem}[Query-local rigidity and exact residual-leaf law]
\label{thm:exact-residual-law}
Under Assumptions~\ref{assump:det-local} and
\ref{assump:hereditary-dag}, for every finite
\(S\subseteq\Cn(B)\) and every \(\widetilde B\subseteq B\),
\begin{align}
  q\in\Cn(\widetilde B\cup S)
  &\Longleftrightarrow
  q\in\Cn\bigl(\widetilde B\cup U(q,S)\bigr),
\label{eq:query-local-rigidity}\\
  &\Longleftrightarrow
  D_{\mathrm{exp}}(q,S)\subseteq\widetilde B.
\label{eq:exact-residual-event}
\end{align}
Consequently, under independent premise erasures with erasure probability
\(\varepsilon\),
\begin{equation}
  \Pr_{\varepsilon}
  \left[
    q\in\Cn(\widetilde B\cup S)
  \right]
  =
  (1-\varepsilon)^{|D_{\mathrm{exp}}(q,S)|}.
\label{eq:exact-residual-probability}
\end{equation}
\end{theorem}

\begin{assumption}[AND--OR local semantics with well-founded expansion]
\label{assump:and-or-local}
The operator \(\Cn\), with its closure axioms and effectiveness, and the
finite relevant region \(\mathcal V_{\mathrm{rel}}(B)\) are retained from
Assumption~\ref{assump:det-local}, as is the base-irreducibility law
\eqref{eq:base-irreducibility}; the designated-parent map, together with its
well-foundedness clause, is superseded by the admissible parent family below
and the rank condition \eqref{eq:well-founded-rank}.  The designated parent
map is replaced by a computable \emph{admissible parent family}: for every
\(v\in\mathcal V_{\mathrm{rel}}(B)\setminus B\),
\[
  \operatorname{Par}(v)
  \subseteq
  2^{\,\mathcal V_{\mathrm{rel}}(B)\setminus\{v\}}
\]
is a nonempty finite family of finite nonempty sets, and for every finite
\(\Gamma\subseteq\Cn(B)\),
\begin{equation}
  \begin{split}
    v & \in\Cn(\Gamma)
    \Longleftrightarrow
    \Bigl[ v\in\Gamma \;\lor \\
       & \exists P\in\operatorname{Par}(v):\, P\subseteq\Cn(\Gamma) \Bigr],
     v\in\mathcal V_{\mathrm{rel}}(B)\setminus B.
  \end{split}
\label{eq:and-or-local-law}
\end{equation}
Thus a derived object is available when it is stored or when \emph{at least
one} of its admissible parent tuples is fully derivable---an AND over each
tuple and an OR over the family.  Write \(r_v:=|\operatorname{Par}(v)|\),
and call \(v\) \emph{ambiguous} if \(r_v\ge 2\).
Assumption~\ref{assump:det-local} is the special case \(r_v=1\) for all
\(v\), with \(\operatorname{par}(v)\) the unique member of
\(\operatorname{Par}(v)\).

Finally, backward expansion is well-founded: there exists a rank function
\[
  \operatorname{rk}:\mathcal V_{\mathrm{rel}}(B)\longrightarrow\mathbb Z_{\ge0},
  \qquad
  \operatorname{rk}^{-1}(0)=B,
\]
such that
\begin{equation}
  u\in P\in\operatorname{Par}(v)
  \quad\Longrightarrow\quad
  \operatorname{rk}(u)<\operatorname{rk}(v).
\label{eq:well-founded-rank}
\end{equation}
\end{assumption}

\begin{theorem}[Envelope rigidity under non-unique derivations]
\label{thm:envelope-rigidity}
Under Assumption~\ref{assump:and-or-local}, suppose \(q\in\Cn(B)\) and let
\(S\subseteq\Cn(B)\) be finite.  Then:
\begin{enumerate}[label=\textup{(\alph*)}]
\item for every \(\widetilde B\subseteq B\),
\begin{equation}
  q\in\Cn(\widetilde B\cup S)
  \Longleftrightarrow
  \exists\,G\in\mathcal G_S(q,B):
  \ D_{\mathrm{exp}}^{G}(q,S)\subseteq\widetilde B;
\label{eq:envelope-event}
\end{equation}
equivalently, with survival indicators
\(X_b:=\mathbf 1\{b\in\widetilde B\}\), the recovery indicator is the
monotone DNF
\begin{equation}
  \mathbf 1\{q\in\Cn(\widetilde B\cup S)\}
  =
  \bigvee_{G\in\mathcal G_S(q,B)}
  \ \bigwedge_{b\in D_{\mathrm{exp}}^{G}(q,S)} X_b;
\label{eq:envelope-dnf}
\end{equation}
\item under independent premise erasures with erasure probability
\(\varepsilon\),
\begin{align}
  (1-\varepsilon)^{e_*(q,S)}
  &\le
  \Pr_{\varepsilon}\!\left[q\in\Cn(\widetilde B\cup S)\right]
  \nonumber\\
  &\le
  \min\!\left\{
    1,\;
    \sum_{G\in\mathcal G_S(q,B)}(1-\varepsilon)^{e_G(q,S)}
  \right\}
  \nonumber\\
  &\le
  \min\!\left\{
    1,\;
    |\mathcal G_S(q,B)|\,(1-\varepsilon)^{e_*(q,S)}
  \right\};
\label{eq:envelope-probability}
\end{align}
\item if the exposed sets
\(\{D_{\mathrm{exp}}^{G}(q,S)\}_{G\in\mathcal G_S(q,B)}\) are pairwise
disjoint, then exactly
\begin{equation}
  \Pr_{\varepsilon}\!\left[q\in\Cn(\widetilde B\cup S)\right]
  =
  1-\prod_{G\in\mathcal G_S(q,B)}
  \Bigl(1-(1-\varepsilon)^{e_G(q,S)}\Bigr);
\label{eq:envelope-independent}
\end{equation}
\item if \(r_v=1\) for every \(v\), then \(e_*(q,S)=|D_{\mathrm{exp}}(q,S)|\)
and the lower bound in \eqref{eq:envelope-probability} coincides with the
exact law \eqref{eq:exact-residual-probability}.
\end{enumerate}
\end{theorem}

\begin{proof}
(a)\ \emph{Sufficiency within a fixed witness.}
Fix \(G\in\mathcal G_S(q,B)\) and suppose
\(D_{\mathrm{exp}}^{G}(q,S)\subseteq\widetilde B\), yet
\(q\notin\Cn(\widetilde B\cup S)\).  Then \(q\notin S\), so \(q\) is not a
source of \(G\), and by \eqref{eq:and-or-local-law} applied with
\(P=N_G^{-}(q)\) some \(u_1\in N_G^{-}(q)\) satisfies
\(u_1\notin\Cn(\widetilde B\cup S)\).  Iterating, any unavailable non-source
vertex has an unavailable in-neighbor in \(G\); unavailable vertices cannot
belong to \(S\), and along edges the rank \eqref{eq:well-founded-rank}
strictly decreases, so the trace terminates at an unavailable source
\(b\).  A source in \(S\) would be available, and an unavailable base source
satisfies \(b\notin\widetilde B\cup S\) by \eqref{eq:base-irreducibility}.
Hence \(b\in B\setminus\widetilde B\), and the traced path
\(b\leadsto q\) meets no vertex of \(U_G(q,S)\), so
\(b\in D_{\mathrm{exp}}^{G}(q,S)\), contradicting
\(D_{\mathrm{exp}}^{G}(q,S)\subseteq\widetilde B\).

\emph{Extraction of a surviving witness.}
Conversely, suppose \(q\in\Cn(\widetilde B\cup S)\) and set
\(A:=\Cn(\widetilde B\cup S)\cap\mathcal V_{\mathrm{rel}}(B)\).  Expand
backward from \(q\): at a visited vertex \(v\in A\), stop if
\(v\in B\cup S\); otherwise \(v\notin\widetilde B\cup S\), and
\eqref{eq:and-or-local-law} with \(\Gamma=\widetilde B\cup S\) yields
\(P_v\in\operatorname{Par}(v)\) with \(P_v\subseteq\Cn(\widetilde B\cup S)\),
hence \(P_v\subseteq A\); continue from each \(u\in P_v\).  Edges strictly
decrease the rank, so the expansion is finite and acyclic, and every visited
vertex reaches \(q\) by construction.  The resulting \(G\) belongs to
\(\mathcal G_S(q,B)\).  Let \(b\in D_{\mathrm{exp}}^{G}(q,S)\) with a
cache-free path \(P:b\leadsto q\).  Walking backward along \(P\) from \(q\),
every internal vertex of \(P\) is a non-source of \(G\), hence was expanded
and lies in \(A\), and its predecessor on \(P\) lies in its chosen parent
tuple, hence also in \(A\).  Therefore \(b\in A\cap B\), so
\eqref{eq:base-irreducibility} gives \(b\in\widetilde B\cup S\); but
\(b\notin S\) since \(P\) is cache-free.  Thus
\(b\in\widetilde B\), proving
\(D_{\mathrm{exp}}^{G}(q,S)\subseteq\widetilde B\).  This completes
\eqref{eq:envelope-event}, and \eqref{eq:envelope-dnf} is a restatement.

(b)\ By \eqref{eq:envelope-event}, recovery is the union of the events
\(\{D_{\mathrm{exp}}^{G}(q,S)\subseteq\widetilde B\}\).  Each exposed set is
a deterministic subset of the distinct base premises, so each event has
probability \((1-\varepsilon)^{e_G(q,S)}\) exactly as in
Theorem~\ref{thm:exact-residual-law}.  The lower bound is the largest single
event probability, and the upper bounds are the union bound and
\(\sum_G(1-\varepsilon)^{e_G(q,S)}\le|\mathcal G_S(q,B)|\,
(1-\varepsilon)^{e_*(q,S)}\).

(c)\ Pairwise disjoint exposed sets make the survival events independent,
and the union probability evaluates to \eqref{eq:envelope-independent}.

(d)\ If \(r_v=1\) for all \(v\), write \(G_q\) for the unique \(B\)-grounded
witness.  Every \(S\)-relative candidate is obtained from \(G_q\) by deleting
the in-edges of a subset of the \(S\)-vertices and pruning vertices that no
longer reach \(q\): expansion is forced except that one may stop at
\(S\)-vertices.  Deleting fewer in-edges can only add directed paths, so the
exposed set is minimized by the full truncation \(T_q\), which deletes the
in-edges of every vertex in \(S\cap V(G_q)\).  The cache-free paths of
\(G_q\) are exactly the paths of \(T_q\) from base leaves to \(q\), so
\(e_*(q,S)=e_{T_q}(q,S)=|D_{\mathrm{exp}}(q,S)|\) with
\(D_{\mathrm{exp}}(q,S)\) as in
Definition~\ref{def:query-local-residual}, and the lower bound of
\eqref{eq:envelope-probability} equals
\eqref{eq:exact-residual-probability}.
\end{proof}

\end{document}
